\section{对称性怎样产生守恒律}
\label{sec:12-Real-Analysis-Support/10-Essential-Structures/01}

一个 N 体引力模拟跑了十万步，行星的位置已经和参考解差出小半圈轨道，可打印出来的总能量仍在初值附近小幅摆动，只有末几位有效数字在动。位置早就不可信，能量却还可信。同一份浮点运算，为什么两个量的命运差这么远？这个反差提示：某些量的稳定性来自比逐点精度更硬的来源。Noether 定理指出那个来源是对称性。

\subsection{守恒量由对称群指定，不靠事后凑}

设位形空间上的轨道为 \(q(t)\)，作用量为 \(S[q]=\int_{t_0}^{t_1}L(q,\dot q,t)\,dt\)。给定单参数变换族 \(\phi_s\)（\(s\in\R\)，\(\phi_0\) 是恒等映射，\(\phi_s\) 对 \(s\) 光滑），若对每条容许轨道都有 \(S[\phi_s\circ q]=S[q]\)，就称 \(\phi_s\) 是作用量的连续对称性。

\begin{theorem}
（Noether，1918）设 \(L\) 对各变元二阶连续可微，单参数光滑变换族 \(\phi_s\) 保持作用量不变，且不移动时间变量。记生成元 \(\xi=\frac{d}{ds}\big|_{s=0}\phi_s(q)\)。则沿任何满足 Euler--Lagrange 方程的轨道，
\[
Q=\sum_i \frac{\partial L}{\partial \dot q_i}\,\xi_i
\]
关于 \(t\) 是常数。
\end{theorem}

证明只有三步。把不变性对 \(s\) 在 \(s=0\) 处求导，得
\[
\begin{aligned}
0&=\frac{d}{ds}\Big|_{s=0}S[\phi_s\circ q]
=\int_{t_0}^{t_1}\Big(\frac{\partial L}{\partial q}\cdot\xi+\frac{\partial L}{\partial \dot q}\cdot\dot\xi\Big)\,dt\\
&=\int_{t_0}^{t_1}\Big(\frac{\partial L}{\partial q}-\frac{d}{dt}\frac{\partial L}{\partial \dot q}\Big)\cdot\xi\,dt
+\Big[\frac{\partial L}{\partial \dot q}\cdot\xi\Big]_{t_0}^{t_1}.
\end{aligned}
\]
第二步是分部积分。轨道满足 Euler--Lagrange 方程时体积分整项为零，只剩边界项；由于 \(t_0,t_1\) 可以任取，\(Q\) 在任意两点取值相同。

条件对账值得停一下。这里用到的是：\(\phi_s\) 对 \(s\) 可微，否则生成元 \(\xi\) 没有定义；不变性对每条容许轨道成立，只在解上成立不够，因为求导那一步要允许离开解的变分；轨道满足 Euler--Lagrange 方程，所以守恒是沿解成立的陈述，而非恒等式。若变换同时平移时间，多出的边界项把 \(Q\) 改写成 \(\sum_i\frac{\partial L}{\partial\dot q_i}\dot q_i-L\)，对 \(L=T-V\) 恰好是能量；空间平移的生成元是常向量，\(Q\) 就是动量。

开头那个模拟也有了答案。leapfrog 一类的辛积分器逐步精确地保持相空间的辛形式，反向误差分析给出：它的离散流是某个与真 Hamiltonian 相差 \(O(h^p)\) 的\emph{影子} Hamiltonian 的精确流。影子量沿数值轨道严格守恒，于是观测到的能量在真值附近有界摆动而不单调漂移。位置误差没有这层保护，只能按 Lyapunov 指数放大。

\subsection{在数据这一侧，对应的问题是哪些变换不改变结论}

把物理换成模型，问句几乎不用改。物理问的是哪些变换保持作用量不变；建模问的是哪些变换应当保持结论不变。手写数字整体右移三像素还是同一个数字，这要求预测函数在平移下不变。分子中原子的输入顺序被打乱，势能预测应当不变，而受力预测应当跟着置换重排。

把它写清楚需要两个概念。设群 \(G\) 作用在输入空间上，\(f\) 是模型。若 \(f(g\cdot x)=f(x)\) 对所有 \(g,x\) 成立，称 \(f\) 不变；若输出空间上另有 \(G\) 的表示 \(\rho\)，且 \(f(g\cdot x)=\rho(g)f(x)\)，称 \(f\) 等变。不变是等变取平凡表示的特例。这两个定义与\hyperref[sec:04-Discrete-Mathematics/06-Groups-and-Symmetry/04]{``不变函数与等变映射：对称性怎样进入计算''}一节完全一致，那里从群作用与轨道讲起，这里关心的是它能兑现成什么。

\begin{center}
\begin{tikzpicture}[>=stealth,line width=0.8pt]
\node at (1.3,3.1) {\footnotesize 硬约束：对所有 \(g\) 成立};
\node (x) at (0,2.2) {\(x\)};
\node (gx) at (2.6,2.2) {\(g\cdot x\)};
\node (fx) at (0,0.6) {\(f(x)\)};
\node (rfx) at (2.6,0.6) {\(\rho(g)f(x)\)};
\draw[->] (x) -- (gx) node[midway,above] {\footnotesize \(g\)};
\draw[->] (x) -- (fx) node[midway,left] {\footnotesize \(f\)};
\draw[->] (gx) -- (rfx) node[midway,right] {\footnotesize \(f\)};
\draw[->] (fx) -- (rfx) node[midway,below] {\footnotesize \(\rho(g)\)};
\begin{scope}[shift={(5.4,0)}]
\node at (1.6,3.1) {\footnotesize 增强：只在抽到的 \(g\) 上成立};
\draw[dashed,line width=0.6pt] (1.6,1.4) ellipse (1.6 and 0.85);
\foreach \a in {25,100,165,255,320} {
  \fill ({1.6+1.6*cos(\a)},{1.4+0.85*sin(\a)}) circle (2pt);
}
\fill (1.6,1.4) circle (1.4pt);
\node[below] at (1.6,1.32) {\footnotesize \(x\)};
\node at (1.6,0.15) {\footnotesize 轨道 \(\set{g\cdot x: g\in G}\)};
\end{scope}
\end{tikzpicture}
\end{center}

\emph{左边的交换图对每一个 \(g\) 都闭合；右边只有实心点处被损失管到，虚线其余部分不受约束。}

\subsection{硬约束写进函数类，增强只写进损失}

硬约束的做法是把等变性做成结构限制，让函数类里根本不含违反它的元素。卷积是最干净的例子：在 \(\Z^d\) 上记平移算子为 \(T_a\)，则线性映射 \(A\) 与所有 \(T_a\) 交换，当且仅当 \(A\) 是与某个核 \(w\) 的卷积。证明骨架是把 \(A\) 作用在单位脉冲上取出 \(w=A\delta_0\)，交换性迫使 \(A\delta_j=T_j w\)，再按线性延拓到一般输入。换句话说，权重共享是平移等变对线性层的唯一解，工程上没有别的选择余地。\hyperref[sec:14-Deep-Learning/04-Convolutional-Networks/01]{``卷积把平移结构写进线性映射''}一节展开了这条刻画在网络中的后果。

一般紧群有一个通用构造。

\begin{proposition}
设紧群 \(G\) 带归一化 Haar 测度 \(\mu\)，\(\rho\) 是输出空间上的表示。对任意映射 \(f\) 定义群平均
\[
\bar f(x)=\int_G \rho(g)^{-1} f(g\cdot x)\,d\mu(g).
\]
则 \(\bar f\) 等变，且 \(f\mapsto\bar f\) 幂等，其值域恰为等变映射构成的子空间。
\end{proposition}

幂等性来自 Haar 测度的平移不变性：对 \(\bar f\) 再做一次平均，被积函数已经与 \(g\) 无关。麻烦也写在定义里：每次前向都要对整个群积分，群大时只能近似，近似误差直接变成等变性的误差。

增强走的是另一条路。经验风险改写成
\[
\frac1n\sum_{i=1}^n \E_{g\sim\mu}\big[\ell\big(f(g\cdot x_i),\,y_i\big)\big],
\]
函数类一点没动，动的是目标。对平方损失有一个可以逐项核对的恒等式：
\[
\E_{g}\big(f(g\cdot x)-y\big)^2=\big(\E_{g}f(g\cdot x)-y\big)^2+\Var_{g}\big(f(g\cdot x)\big).
\]
右端第一项是拿轨道平均去拟合标签，第二项是一个惩罚，它为零当且仅当 \(f\) 在这条轨道上几乎处处取常值。所以增强在平方损失下的作用是明确的：把不变性变成一项系数固定为 \(1\) 的正则。

两种兑现的强度差别就落在这里。回看上图右侧：惩罚只在训练样本的轨道上生效，只覆盖 \(\mu\) 的支撑，而且它是惩罚不是约束——当拟合项和它冲突时，优化器会拿不变性去换训练误差。左侧的交换图则对每个 \(x\)、每个 \(g\) 恒等成立，与数据和优化都无关。把增强称作``软化版本''是准确的，把它当作等变性的等价替代则不准确。

\subsection{破缺不让守恒律失效，只让它的适用范围缩小}

Noether 定理要求精确不变性，但破缺的后果是可以定量的。若 \(L\) 显含时间，直接计算给出 \(\frac{dE}{dt}=-\frac{\partial L}{\partial t}\)：这是等式，破缺多大，漂移就多快，慢变外场对应缓慢漂移而非突然失控。

\begin{remark}
数据一侧的对应说法要弱一些。有限画幅上的平移等变只是近似：像素移出边界后信息丢失，带步长的下采样让等变只对步长整数倍的平移成立，位置编码则直接把绝对坐标写进输入。这些是对具体架构的观察，不是定理；要变成定理，得先把边界条件与采样格点写进模型的定义域。
\end{remark}

\subsection{把定理与启发分开放}

这一节里够得上定理的有三条：Noether 定理，平移交换算子的卷积刻画，紧群上群平均的幂等性。够得上恒等式的有一条：平方损失下增强等于轨道平均拟合加方差惩罚。剩下的``增强能带来不变性''属于启发，它在训练分布之外没有保证。

由此可以留下一条操作性的问法：先确定变换群是什么、它是否也作用在输出上，再确定这份对称性是恒等成立还是只在样本上兑现。对称性削减的是模型必须表达的自由度。下一节从另一头问同样的削减——当数据本身的自由度远低于它所在空间的维数时，凭什么相信这件事会发生。
